% =============================================================================
% WheelsOn - Documento del Proyecto de Investigación y Desarrollo (I+D)
% Bases de Datos - Bases de datos orientadas a grafos
% Compilar con XeLaTeX o Tectonic:  tectonic main.tex
% =============================================================================
\documentclass[11pt, a4paper]{article}

% ----------------------------------------------------------------- paquetes
\usepackage{fontspec}
\usepackage[spanish, es-noshorthands, es-tabla]{babel}
\usepackage[margin=2.5cm, headheight=16pt]{geometry}
\usepackage[table]{xcolor}
\usepackage{titlesec}
\usepackage{fancyhdr}
\usepackage{graphicx}
\usepackage{float}
\usepackage[section]{placeins}
\usepackage{caption}
\usepackage{listings}
\usepackage{enumitem}
\usepackage{booktabs}
\usepackage{array}
\usepackage{tabularx}
\usepackage{xltabular}
\usepackage{ragged2e}
\usepackage[most]{tcolorbox}
\usepackage{tikz}
\usetikzlibrary{arrows.meta, positioning, calc}
\usepackage{microtype}
\usepackage{needspace}
\usepackage[hidelinks]{hyperref}

% ----------------------------------------------------------------- fuentes
\setmainfont{texgyrepagella}[
  Extension = .otf,
  UprightFont = *-regular,
  BoldFont = *-bold,
  ItalicFont = *-italic,
  BoldItalicFont = *-bolditalic]
\setmonofont{Inconsolatazi4}[
  Extension = .otf,
  UprightFont = *-Regular,
  BoldFont = *-Bold,
  Scale = 0.92]

% ----------------------------------------------------------------- colores
\definecolor{azul}      {RGB}{20, 48, 89}     % azul marino principal
\definecolor{azulmedio} {RGB}{41, 88, 150}
\definecolor{azulclaro} {RGB}{152, 180, 216}
\definecolor{azulfondo} {RGB}{236, 242, 250}
\definecolor{azulfila}  {RGB}{245, 248, 252}
\definecolor{tinta}     {RGB}{33, 37, 43}
\definecolor{gris}      {RGB}{100, 108, 120}
\definecolor{borde}     {RGB}{205, 216, 230}
\definecolor{codfondo}  {RGB}{16, 27, 45}
\definecolor{codtexto}  {RGB}{226, 233, 242}
\definecolor{codclave}  {RGB}{122, 178, 255}
\definecolor{codcoment} {RGB}{128, 146, 170}
\definecolor{codtexto2} {RGB}{158, 214, 172}

\color{tinta}

% ----------------------------------------------------------------- encabezado
\pagestyle{fancy}
\fancyhf{}
\fancyhead[L]{\textcolor{azul}{\small\textbf{WheelsOn}}}
\fancyhead[R]{\textcolor{azul}{\small\textbf{Bases de Datos --- Grafos con Neo4j}}}
\fancyfoot[C]{\textcolor{gris}{\small\thepage}}
\renewcommand{\headrulewidth}{0.8pt}
\renewcommand{\headrule}{\hbox to\headwidth{\color{azul}\leaders\hrule height \headrulewidth\hfill}}

% ----------------------------------------------------------------- títulos
\titleformat{\section}
  {\Large\bfseries\color{azul}}{\thesection.}{0.6em}{}
  [\vspace{-0.5em}\textcolor{azul}{\rule{\linewidth}{0.8pt}}]
\titleformat{\subsection}
  {\large\bfseries\color{tinta}}{\thesubsection.}{0.5em}{}
\titleformat{\subsubsection}
  {\normalsize\bfseries\color{azulmedio}}{}{0em}{}
\pretocmd{\section}{\needspace{10\baselineskip}}{}{}
\pretocmd{\subsection}{\needspace{6\baselineskip}}{}{}
\titlespacing*{\section}{0pt}{1.6em}{0.9em}
\titlespacing*{\subsection}{0pt}{1.2em}{0.5em}

\setlength{\parindent}{1.2em}
\setlength{\parskip}{0.45em}
\linespread{1.08}
\setlist{itemsep=0.2em, topsep=0.3em}

\captionsetup{font=small, labelfont={bf,color=azul}, labelsep=period}
\captionsetup[table]{position=top, skip=6pt}

% ----------------------------------------------------------------- tablas
\newcolumntype{L}{>{\RaggedRight\arraybackslash}X}
\newcommand{\cabecera}{\rowcolor{azul}}
\newcommand{\h}[1]{\textcolor{white}{\textbf{#1}}}
\renewcommand{\arraystretch}{1.25}

% ----------------------------------------------------------------- recuadros
\newtcolorbox{nota}[1]{
  enhanced, breakable,
  colback=azulfondo, colframe=azul,
  boxrule=0.6pt, arc=2pt,
  left=8pt, right=8pt, top=5pt, bottom=6pt,
  fonttitle=\bfseries\small, coltitle=white, colbacktitle=azul,
  title={#1}, before skip=10pt, after skip=10pt}

% ----------------------------------------------------------------- código
\lstdefinelanguage{Cypher}{
  morekeywords={MATCH, OPTIONAL, WHERE, RETURN, WITH, ORDER, BY, LIMIT, DESC, ASC,
    AND, OR, NOT, EXISTS, CREATE, CONSTRAINT, INDEX, IF, FOR, REQUIRE, IS, UNIQUE,
    ON, DISTINCT, AS, CASE, WHEN, THEN, END, ELSE, IN, COUNT, DETACH, DELETE, NULL,
    shortestPath, count, collect, size, coalesce, toLower, all, any, type, nodes,
    relationships, length, labels},
  sensitive=true,
  morecomment=[l]{//},
  morestring=[b]', morestring=[b]"
}
\lstdefinestyle{oscuro}{
  basicstyle=\ttfamily\small\color{codtexto},
  backgroundcolor=\color{codfondo},
  keywordstyle=\bfseries\color{codclave},
  commentstyle=\itshape\color{codcoment},
  stringstyle=\color{codtexto2},
  frame=single, framerule=0pt, rulecolor=\color{codfondo},
  framexleftmargin=8pt, framexrightmargin=8pt, framextopmargin=6pt, framexbottommargin=6pt,
  xleftmargin=8pt, xrightmargin=8pt,
  breaklines=true, breakatwhitespace=true, columns=fullflexible, keepspaces=true,
  showstringspaces=false, tabsize=2,
  aboveskip=12pt, belowskip=10pt,
  extendedchars=true, inputencoding=utf8,
  literate={á}{{\'a}}1 {é}{{\'e}}1 {í}{{\'i}}1 {ó}{{\'o}}1 {ú}{{\'u}}1 {ñ}{{\~n}}1
           {Á}{{\'A}}1 {É}{{\'E}}1 {Í}{{\'I}}1 {Ó}{{\'O}}1 {Ú}{{\'U}}1
}
\lstset{style=oscuro, language=Cypher}
\lstdefinestyle{anexo}{style=oscuro, language=Cypher,
  basicstyle=\ttfamily\footnotesize\color{tinta},
  backgroundcolor=\color{azulfila}, rulecolor=\color{borde}, framerule=0.4pt,
  keywordstyle=\bfseries\color{azulmedio}, commentstyle=\itshape\color{gris},
  stringstyle=\color{teal!70!black}, breakatwhitespace=false}
\lstdefinestyle{sql}{style=oscuro, language=SQL,
  morekeywords={RECURSIVE, ARRAY, ANY, UNION, ALL, REFERENCES, PRIMARY, KEY, FOREIGN}}

% ----------------------------------------------------------------- utilidades
\newcommand{\pend}[1]{{\color{gris}\mdseries\itshape #1}}
\newcommand{\cod}[1]{\texttt{\detokenize{#1}}}
\newcommand{\rel}[1]{\texttt{\textbf{\detokenize{#1}}}}
\newcommand{\nodo}[1]{\texttt{:#1}}
\newcommand{\falta}[1]{\colorbox{red!12}{\textcolor{red!70!black}{\textbf{Falta:} #1}}}
\newif\ifcapturafija
\newcommand{\captura}[4][0.99\linewidth]{%
  \ifcapturafija\begin{figure}[H]\else\begin{figure}[htbp]\fi
    \centering
    {\setlength{\fboxsep}{0pt}\setlength{\fboxrule}{0.6pt}%
     \fcolorbox{borde}{white}{\includegraphics[width=#1]{img/#2}}}
    \caption{#3}\label{fig:#4}
  \end{figure}}

% ----------------------------------------------------------------- tikz
\tikzset{
  ent/.style={draw=azul, fill=azulfondo, rounded corners=5pt, line width=0.9pt,
    minimum width=2.7cm, minimum height=0.95cm, font=\bfseries\color{azul}},
  rel/.style={-{Stealth[length=2.4mm]}, line width=0.8pt, color=azulmedio},
  et/.style={font=\scriptsize\ttfamily\bfseries, fill=white, inner sep=1.6pt,
    text=tinta, align=center},
}

% =============================================================================
\renewcommand{\topfraction}{0.85}
\renewcommand{\bottomfraction}{0.7}
\renewcommand{\textfraction}{0.12}
\renewcommand{\floatpagefraction}{0.75}
\setcounter{topnumber}{3}
\setcounter{totalnumber}{4}

\begin{document}

% ----------------------------------------------------------------- portada
\begin{titlepage}
  \newgeometry{margin=2.6cm}
  \pagecolor{azul}
  \color{white}
  \setlength{\parindent}{0pt}
  {\small\bfseries\scshape Proyecto de Investigación y Desarrollo (I+D) --- Bases de Datos}\\[0.25cm]
  \textcolor{azulclaro}{\rule{\linewidth}{0.5pt}}

  \vspace{3.2cm}
  {\fontsize{50}{54}\selectfont\bfseries WheelsOn}\\[0.5cm]
  {\Large Una red profesional sobre una base de datos de grafos}

  \vspace{0.8cm}
  \textcolor{azulclaro}{\rule{\linewidth}{0.5pt}}\\[0.3cm]
  \begin{minipage}{0.8\linewidth}
    \large\color{azulfondo}
    Diseño e implementación de una aplicación basada en una base de datos de grafos
    para la gestión y el análisis de relaciones en una red profesional,
    construida con Neo4j y Django.
  \end{minipage}

  \vfill
  \begin{tabular}{@{}l@{\hspace{1.2cm}}l@{}}
    Dubal Antonio Aguilar Torres & 200207378 \\
    Carlos Miguel Calabria Patiño & 200209929 \\
    Angelo David de Leon Caneda & 200205958 \\
    Mateo José Cisneros Galeano & 200205979 \\
    Luis Angel Cuellar Ballesteros & 200163979 \\
  \end{tabular}

  \vspace{1.6cm}
  \begin{center}
    {\bfseries Docente:}\\[0.15cm]
    Wilson Manuel Estrada Bernal\\[1.2cm]
    Ingeniería de Sistemas\\
    Universidad del Norte\\
    Barranquilla, Colombia\\[0.3cm]
    Noviembre de 2026
  \end{center}
  \restoregeometry
\end{titlepage}
\nopagecolor
\color{tinta}

\setcounter{tocdepth}{2}
\begingroup\small\setlength{\parskip}{0.1em}\tableofcontents\endgroup
\newpage

% =============================================================================
\section{Introducción}
% =============================================================================

Desarrollamos WheelsOn, una red profesional en la que cada persona tiene un perfil, se
conecta con otras, declara sus habilidades, trabaja en empresas, participa en proyectos
y se postula a ofertas de empleo. En una red así importa más cómo se
relacionan los datos que cada dato por separado: quién conoce a quién, qué habilidades comparten dos personas o por qué
contactos se puede llegar a una empresa.

Por eso guardamos toda la información de la red en una base de datos de grafos, Neo4j,
y construimos la aplicación web en Python con Django.

En este documento explicamos primero qué es una base de datos de grafos
(Sección~\ref{sec:grafos}) y presentamos la aplicación (Sección~\ref{sec:app}). Luego
describimos el modelo que diseñamos para WheelsOn (Sección~\ref{sec:modelo}), los
scripts con los que creamos la base y cargamos los datos (Sección~\ref{sec:scripts}) y
las consultas con las que respondemos las preguntas del enunciado
(Sección~\ref{sec:consultas}). Por último, a partir de lo que aprendimos al construir
WheelsOn, analizamos en qué escenarios el modelo de grafos supera al relacional y en
cuáles no (Sección~\ref{sec:investigacion}).

% =============================================================================
\section{Bases de datos de grafos}\label{sec:grafos}
% =============================================================================

\subsection{La información como una red}

Una base de datos de grafos guarda la información como una red de puntos unidos por
líneas. Los puntos se llaman \textbf{nodos} y representan cosas: una persona, una
empresa, una habilidad. Las líneas se llaman \textbf{relaciones} y dicen cómo se
vinculan esas cosas, por ejemplo, ``Valentina trabaja en Andina Logística''.

Neo4j usa el \emph{modelo de grafo de propiedades}, que tiene cuatro piezas:

\begin{table}[htbp]
\centering
\caption{Piezas del modelo de grafo de propiedades.}
\rowcolors{2}{azulfila}{white}
\begin{tabularx}{\linewidth}{@{}l L l@{}}
\cabecera \h{Pieza} & \h{Qué es} & \h{Ejemplo en WheelsOn} \\
Nodo      & Una entidad.                                   & Valentina Ríos \\
Etiqueta  & El tipo de un nodo.                            & \nodo{Usuario}, \nodo{Empresa} \\
Relación  & Un vínculo entre dos nodos, con tipo y dirección. & \rel{TRABAJA_EN} \\
Propiedad & Un dato guardado en un nodo o en una relación. & \cod{nombre}, \cod{cargo} \\
\bottomrule
\end{tabularx}
\end{table}

A diferencia de otros modelos, aquí las relaciones también guardan datos. En
la Figura~\ref{fig:ejemplo}, el cargo y la fecha de inicio están en la relación
\rel{TRABAJA_EN}, porque son datos del vínculo entre Valentina y la empresa.

\begin{figure}[htbp]
\centering
\begin{tikzpicture}
  \node[ent, align=center] (u) at (0,0) {Valentina Ríos\\[-1pt]{\scriptsize\mdseries\ttfamily :Usuario}};
  \node[ent, align=center] (e) at (8.2,0) {Andina Logística\\[-1pt]{\scriptsize\mdseries\ttfamily :Empresa}};
  \draw[rel] (u) -- node[et, above=2pt] {TRABAJA\_EN}
                    node[below=3pt, font=\scriptsize\ttfamily, text=gris, align=center]
                    {cargo: "Directora de Tecnología"\\ fecha\_inicio: "2021-03-01"\\ actual: true} (e);
\end{tikzpicture}
\caption{Dos nodos y una relación con propiedades, tomados de los datos de demostración.}
\label{fig:ejemplo}
\end{figure}

\subsection{Por qué recorrer un grafo es rápido}

En una base relacional, los vínculos entre registros se reconstruyen en cada consulta:
el motor busca, tabla por tabla, las filas cuyas claves coinciden (la operación
\texttt{JOIN}). Cada nivel más de profundidad en la pregunta agrega otro \texttt{JOIN}.

En Neo4j, cada nodo guarda referencias directas a sus relaciones. Pasar de un nodo a
sus vecinos es seguir esas referencias, sin buscar en una tabla. Por eso el costo de
un recorrido depende de cuántos nodos visita, no de cuántos datos hay en total en la
base. Esta propiedad se conoce como \emph{adyacencia sin índice}.

\subsection{Cypher, el lenguaje de consultas}

Neo4j se consulta con Cypher. Una consulta describe la forma que deben tener los datos
con un dibujo hecho de texto: los nodos van entre paréntesis y las relaciones entre
corchetes, con una flecha que indica la dirección. Esta consulta devuelve quiénes
trabajan en Andina Logística y con qué cargo:

\begin{lstlisting}
MATCH (u:Usuario)-[t:TRABAJA_EN]->(e:Empresa {nombre: 'Andina Logística'})
RETURN u.nombre, t.cargo
\end{lstlisting}

\cod{MATCH} describe el patrón que se busca y \cod{RETURN} lo que se quiere obtener.
Además, en Cypher la cantidad de saltos puede ser variable:
\cod{-[:CONECTA_CON*..6]-} significa ``por una cadena de entre uno y seis contactos''.
Con eso se escriben en una línea preguntas como la ruta más corta entre dos personas.

% =============================================================================
\section{La aplicación}\label{sec:app}
% =============================================================================

Construimos WheelsOn como una aplicación web en la que cada pantalla lee y escribe
directamente en Neo4j.

\begin{table}[htbp]
\centering
\caption{Tecnologías utilizadas.}
\rowcolors{2}{azulfila}{white}
\begin{tabularx}{\linewidth}{@{}l L@{}}
\cabecera \h{Parte} & \h{Tecnología} \\
Base de datos      & Neo4j 5.26, en un contenedor de Docker. Django usa además un archivo SQLite local solo para las sesiones de inicio de sesión; ningún dato de la red se guarda ahí. \\
Aplicación web     & Python 3.12 y Django 6. \\
Conexión con Neo4j & \emph{neomodel}, que convierte nodos en objetos de Python, y consultas Cypher escritas a mano para todo lo que recorre el grafo. \\
Interfaz           & HTML con Tailwind y Alpine.js; el grafo interactivo usa vis-network y las gráficas, Chart.js. \\
\bottomrule
\end{tabularx}
\end{table}

\captura[0.78\linewidth]{inicio.png}{Pantalla de inicio de WheelsOn con la sesión de Valentina Ríos, una de las
personas de los datos de demostración.}{inicio}

En la Tabla~\ref{tab:funcionalidades} ubicamos cada funcionalidad del enunciado dentro
de la aplicación.

\begin{table}[htbp]
\centering
\caption{Funcionalidades pedidas y dónde se encuentran.}\label{tab:funcionalidades}
\rowcolors{2}{azulfila}{white}
\small
\begin{tabularx}{\linewidth}{@{}>{\RaggedRight\arraybackslash}p{5.6cm} L@{}}
\cabecera \h{Funcionalidad} & \h{En la aplicación} \\
CRUD de usuarios & Registro, perfil, edición y eliminación de la cuenta. \\
Conexiones entre usuarios & Solicitudes de conexión que la otra persona acepta o rechaza. \\
Vincular usuarios con habilidades, empresas y proyectos & Desde el perfil, con nivel, cargo o rol. Los empleos y la entrada a proyectos los aprueba la empresa o el líder. \\
Empresas y sus ofertas & Páginas de Empresas y Empleos. \\
Habilidades requeridas por una oferta & Requisitos de cada oferta, con nivel mínimo y prioridad. \\
Postularse a una oferta & Botón \emph{Postularme}. La empresa mueve la postulación por etapas hasta aceptarla o rechazarla. \\
Caminos y conexiones & Consultas Cypher: rutas entre personas, hacia empresas y hacia ofertas (Sección~\ref{sec:consultas}). \\
Recomendaciones & Sugerencias de contactos, empleos, colaboradores para proyectos y talento para empresas. \\
Panel de analítica & Influencia, comunidades y habilidades más demandadas (Sección~\ref{sec:analitica}). \\
\bottomrule
\end{tabularx}
\end{table}

% =============================================================================
\section{Modelo utilizado}\label{sec:modelo}
% =============================================================================

\subsection{Visión general}

Nuestro modelo tiene 8 tipos de nodo y 23 tipos de relación (Figura~\ref{fig:modelo}).
Conservamos los nombres que propone el enunciado y agregamos los nodos y relaciones que
hicieron falta para que la aplicación funcionara como una red real: publicaciones,
avances de proyecto, notificaciones, permisos, alianzas entre empresas y solicitudes
pendientes.

\begin{figure}[htbp]
\centering
\begin{tikzpicture}[x=0.85cm, y=0.9cm, ent/.append style={minimum width=2.4cm}]
  \node[ent] (h)  at (-4, 5.2)   {Habilidad};
  \node[ent] (o)  at ( 4, 5.2)   {Oferta};
  \node[ent] (n)  at (-8, 2.2)   {Notificacion};
  \node[ent] (u)  at ( 0, 0)     {Usuario};
  \node[ent] (e)  at ( 8, 0)     {Empresa};
  \node[ent] (a)  at (-8,-2.8)   {Avance};
  \node[ent] (p)  at (-4,-6)     {Proyecto};
  \node[ent] (pb) at ( 4,-6)     {Publicacion};

  \draw[rel] (o) -- node[et] {REQUIERE} (h);
  \draw[rel] (u) -- node[et, pos=0.55] {TIENE\_HABILIDAD} (h);
  \draw[rel] (u) -- node[et, pos=0.55] {POSTULA\_A} (o);
  \draw[rel] (e) -- node[et] {PUBLICA} (o);
  \draw[rel] (u) to[out=72, in=108, looseness=7] node[et, above=1pt] {CONECTA\_CON\\\pend{SOLICITA\_CONEXION}} (u);
  \draw[rel] (u) -- node[et] {TRABAJA\_EN\\ADMINISTRA\\\pend{SOLICITA\_EMPLEO}\\\pend{INVITADO\_A\_TRABAJAR}\\\pend{INVITADO\_A\_ADMINISTRAR}} (e);
  \draw[rel] (e) to[out=-62, in=-118, looseness=6] node[et, below=1pt] {ALIADA\_DE\\\pend{SOLICITA\_ALIANZA}} (e);
  \draw[rel] (u) -- node[et] {RECIBE} (n);
  \draw[rel] (u) -- node[et, pos=0.5] {REGISTRA} (a);
  \draw[rel] (p) -- node[et] {TIENE\_AVANCE} (a);
  \draw[rel] (u) -- node[et, pos=0.55] {PARTICIPA\_EN\\LIDERA\\\pend{SOLICITA\_UNIRSE}\\\pend{INVITADO\_A}} (p);
  \draw[rel] (u) -- node[et, pos=0.55] {AUTOR\_DE\\REDACTA\\PERCIBE} (pb);
  \draw[rel] (e) -- node[et, pos=0.45] {AUTOR\_DE} (pb);
\end{tikzpicture}
\caption{Modelo de WheelsOn: los 8 tipos de nodo y los 23 tipos de relación. En gris,
las solicitudes e invitaciones pendientes, que se reemplazan por la relación
definitiva cuando se aceptan (Tabla~\ref{tab:pendientes}).}
\label{fig:modelo}
\end{figure}

\subsection{Nodos}

A cada nodo le dimos un identificador propio (\cod{uid}, \cod{eid}, \cod{oid}, etc.)
que no se repite. Es el que aparece en las direcciones de la aplicación, por ejemplo
\cod{/usuarios/usr_valentina16/}.

\begin{table}[htbp]
\centering
\caption{Tipos de nodo.}
\rowcolors{2}{azulfila}{white}
\small
\begin{tabularx}{\linewidth}{@{}l >{\RaggedRight\arraybackslash}p{3.8cm} L@{}}
\cabecera \h{Etiqueta} & \h{Representa} & \h{Propiedades principales} \\
\nodo{Usuario}      & Una persona registrada.               & \cod{uid}, \cod{nombre}, \cod{apellido}, \cod{email}, \cod{titulo_profesional}, \cod{ubicacion}, \cod{rol} \\
\nodo{Empresa}      & Una organización.                     & \cod{eid}, \cod{nombre}, \cod{industria}, \cod{tamano}, \cod{ubicacion} \\
\nodo{Habilidad}    & Un conocimiento o competencia.        & \cod{hid}, \cod{nombre}, \cod{categoria} (técnica, blanda o idioma) \\
\nodo{Proyecto}     & Un trabajo o iniciativa.              & \cod{pid}, \cod{nombre}, \cod{estado}, \cod{tecnologias} (lista), fechas \\
\nodo{Oferta}       & Una vacante laboral.                  & \cod{oid}, \cod{titulo}, \cod{modalidad}, \cod{salario_min}, \cod{salario_max}, \cod{cupos}, \cod{estado} \\
\nodo{Publicacion}  & Lo que alguien comparte en el inicio. & \cod{pubid}, \cod{texto}, \cod{fecha} \\
\nodo{Avance}       & Una novedad de un proyecto.           & \cod{aid}, \cod{texto}, \cod{estado}, \cod{fecha} \\
\nodo{Notificacion} & Un aviso para una persona.            & \cod{nid}, \cod{titulo}, \cod{detalle}, \cod{fecha} \\
\bottomrule
\end{tabularx}
\end{table}

\FloatBarrier
\subsection{Relaciones}

{\small
\rowcolors{2}{azulfila}{white}
\begin{xltabular}{\linewidth}{@{}l >{\RaggedRight\arraybackslash}p{3.4cm} >{\RaggedRight\arraybackslash}p{3.4cm} L@{}}
\caption{Relaciones del modelo.}\label{tab:relaciones}\\
\cabecera \h{Relación} & \h{De $\rightarrow$ a} & \h{Propiedades} & \h{Significado} \\
\endfirsthead
\cabecera \h{Relación} & \h{De $\rightarrow$ a} & \h{Propiedades} & \h{Significado} \\
\endhead
\rel{CONECTA_CON}     & Usuario -- Usuario    & \cod{fecha}, \cod{mensaje}, \cod{estado} (siempre \emph{aceptada}) & Contacto confirmado por ambas personas. Hay una sola relación por pareja y las consultas la recorren sin importar la dirección. \\
\rel{TIENE_HABILIDAD} & Usuario $\rightarrow$ Habilidad & \cod{nivel}, \cod{anios_experiencia} & La persona domina esa habilidad. \\
\rel{TRABAJA_EN}      & Usuario $\rightarrow$ Empresa & \cod{cargo}, \cod{fecha_inicio}, \cod{fecha_fin}, \cod{actual} & Empleo actual o anterior. \\
\rel{ADMINISTRA}      & Usuario $\rightarrow$ Empresa & --- & Quién gestiona la empresa y sus ofertas. \\
\rel{PARTICIPA_EN}    & Usuario $\rightarrow$ Proyecto & \cod{rol}, \cod{fecha_inicio}, \cod{activo} & Miembro del equipo de un proyecto. \\
\rel{LIDERA}          & Usuario $\rightarrow$ Proyecto & --- & Quién gestiona el proyecto y su equipo. \\
\rel{PUBLICA}         & Empresa $\rightarrow$ Oferta & \cod{fecha_publicacion} & La empresa ofrece esa vacante. \\
\rel{REQUIERE}        & Oferta $\rightarrow$ Habilidad & \cod{nivel_minimo}, \cod{prioridad} & Requisito de la vacante, obligatorio o deseable. \\
\rel{POSTULA_A}       & Usuario $\rightarrow$ Oferta & \cod{estado}, \cod{fecha}, \cod{fecha_estado}, \cod{carta_presentacion} & Postulación y la etapa en que va. \\
\rel{ALIADA_DE}       & Empresa -- Empresa & \cod{fecha} & Alianza aceptada entre dos empresas. \\
\rel{AUTOR_DE}        & Usuario o Empresa $\rightarrow$ Publicacion & --- & A nombre de quién aparece la publicación. \\
\rel{REDACTA}         & Usuario $\rightarrow$ Publicacion & --- & Qué persona escribió una publicación de empresa. \\
\rel{PERCIBE}         & Usuario $\rightarrow$ Publicacion & \cod{tipo} (buena o mala), \cod{fecha} & Reacción a una publicación. \\
\rel{TIENE_AVANCE}    & Proyecto $\rightarrow$ Avance & --- & Historial del proyecto. \\
\rel{REGISTRA}        & Usuario $\rightarrow$ Avance & --- & Quién registró el avance. \\
\rel{RECIBE}          & Usuario $\rightarrow$ Notificacion & --- & Avisos de la persona. \\
\bottomrule
\end{xltabular}
}

\begin{table}[htbp]
\centering
\caption{Solicitudes e invitaciones pendientes. Al aceptarse, cada una se reemplaza por la relación definitiva.}\label{tab:pendientes}
\rowcolors{2}{azulfila}{white}
\small
\begin{tabularx}{\linewidth}{@{}l l L@{}}
\cabecera \h{Pendiente} & \h{De $\rightarrow$ a} & \h{Al aceptarse se convierte en} \\
\rel{SOLICITA_CONEXION}      & Usuario $\rightarrow$ Usuario  & \rel{CONECTA_CON} \\
\rel{SOLICITA_EMPLEO}        & Usuario $\rightarrow$ Empresa  & \rel{TRABAJA_EN}, cuando la empresa confirma el empleo. \\
\rel{INVITADO_A_TRABAJAR}    & Usuario $\rightarrow$ Empresa  & \rel{TRABAJA_EN}, cuando la persona acepta. \\
\rel{INVITADO_A_ADMINISTRAR} & Usuario $\rightarrow$ Empresa  & \rel{ADMINISTRA}: la administración pasa a esa persona. \\
\rel{SOLICITA_UNIRSE}        & Usuario $\rightarrow$ Proyecto & \rel{PARTICIPA_EN} \\
\rel{INVITADO_A}             & Usuario $\rightarrow$ Proyecto & \rel{PARTICIPA_EN} o \rel{LIDERA}, según la invitación. \\
\rel{SOLICITA_ALIANZA}       & Empresa $\rightarrow$ Empresa  & \rel{ALIADA_DE} \\
\bottomrule
\end{tabularx}
\end{table}

\subsection{Decisiones de diseño}

\subsubsection{Los datos de un vínculo van en la relación}
Pusimos el cargo en \rel{TRABAJA_EN}, la etapa de la postulación en \rel{POSTULA_A} y
la prioridad de un requisito en \rel{REQUIERE}. Si el cargo estuviera en el nodo de la
persona, nadie podría tener dos empleos.

\subsubsection{Lo pendiente vive en relaciones aparte}
Guardamos una solicitud de conexión como una relación \rel{SOLICITA_CONEXION}. Si se
acepta, la borramos y creamos \rel{CONECTA_CON}; si se rechaza, solo la borramos. Así, las rutas, las
recomendaciones y la analítica recorren únicamente \rel{CONECTA_CON} y nunca cuentan
como contacto a alguien que no ha respondido. Aplicamos la misma idea a empleos, proyectos,
administración y alianzas (Tabla~\ref{tab:pendientes}).

\subsubsection{Una postulación avanza por etapas}
\rel{POSTULA_A} pasa por \emph{enviada}, \emph{en revisión} y \emph{entrevista}, hasta
terminar en \emph{aceptada}, \emph{rechazada} o \emph{no seleccionada}. Aceptar a alguien
crea su relación \rel{TRABAJA_EN} con el título del empleo como cargo. Cuando se llenan
los cupos, la oferta se cierra y quienes seguían en proceso quedan como \emph{no
seleccionada}. Los empleos terminados los conservamos con \cod{actual: false} para la
trayectoria de la persona, pero en las rutas solo usamos los actuales.

\subsubsection{Los permisos también están en el grafo}
Decidimos quién puede editar algo con una relación. Para cambiar una oferta, por
ejemplo, la aplicación comprueba que exista este camino (o que la persona sea
administradora de la plataforma):
\begin{center}\cod{(persona)-[:ADMINISTRA]->(empresa)-[:PUBLICA]->(oferta)}\end{center}

\subsubsection{Las notificaciones son nodos}
Lo que está pendiente se deduce de las relaciones de la Tabla~\ref{tab:pendientes}.
Pero cuando alguien rechaza una solicitud, la relación se borra y ya no queda nada de
qué avisar. Por eso guardamos cada aviso como un nodo \nodo{Notificacion} unido a la
persona con \rel{RECIBE}.

\subsubsection{Lo que dejamos como propiedad}
Guardamos las tecnologías de un proyecto como una lista de texto (\cod{tecnologias}), no
como relaciones con nodos \nodo{Habilidad}. Esto simplificó la creación de proyectos,
pero nos obliga a comparar texto cuando buscamos proyectos parecidos (Sección~\ref{sec:proyectos}). Como
mejora, proponemos una relación \cod{USA_TECNOLOGIA} hacia \nodo{Habilidad}.

\captura[0.62\linewidth]{grafo.png}{Grafo interactivo de la aplicación: los datos de demostración tal como
están en Neo4j. Círculos azules: usuarios; cuadrados turquesa: empresas; rombos: habilidades;
hexágonos morados: proyectos; triángulos rojos: ofertas.}{grafo}

\subsection{Restricciones e índices}\label{sec:restricciones}

Las \textbf{restricciones de unicidad} impiden que se repita un valor. Hay once: el
identificador de cada uno de los ocho tipos de nodo, el correo de cada usuario y los
nombres de empresas y habilidades. Los \textbf{índices} aceleran las búsquedas
frecuentes. Hay cinco: el nombre y apellido de los usuarios, la industria de las
empresas, la categoría de las habilidades y el estado de las ofertas y de los
proyectos.

Neo4j no puede exigir por sí mismo el resto de las reglas del negocio, como que una oferta
pertenezca a una sola empresa o que no se acepte a más personas que cupos. Esas las
hicimos cumplir desde la aplicación, dentro de transacciones
(Sección~\ref{sec:no-gana}).

% =============================================================================
\section{Scripts de inicio de la base de datos}\label{sec:scripts}
% =============================================================================

Organizamos los scripts en la carpeta \cod{scripts/} del repositorio.

\begin{table}[htbp]
\centering
\caption{Scripts de la base de datos, en el orden en que se ejecutan.}
\rowcolors{2}{azulfila}{white}
\small
\begin{tabularx}{\linewidth}{@{}l L@{}}
\cabecera \h{Archivo} & \h{Contenido} \\
\cod{constraints.cypher} & Las 11 restricciones y los 5 índices de la Sección~\ref{sec:restricciones}. Usa \cod{IF NOT EXISTS}, así que se puede ejecutar varias veces sin error. Va primero (completo en el Anexo~\ref{anx:inicio}). \\
\cod{init_db.cypher}     & Lo mismo, una sentencia por línea, para pegarlo en Neo4j Browser. \\
\cod{seed_data.cypher}   & Datos base: 15 personas, 3 empresas, 12 habilidades, 3 proyectos y 2 ofertas, con sus relaciones. \\
\cod{datos_demo.cypher}  & La red de demostración: 25 personas, 8 empresas y 13 ofertas más, publicaciones, avances, postulaciones en todas las etapas y solicitudes pendientes. \\
\cod{queries.cypher}     & Las consultas de la Sección~\ref{sec:consultas} (completo en el Anexo~\ref{anx:consultas}). \\
\bottomrule
\end{tabularx}
\end{table}

\subsection{Carga de datos}

Primero creamos los nodos y después las relaciones. Para crear una relación, buscamos
sus dos extremos por su identificador con \cod{MATCH} y los unimos con \cod{CREATE}:

\begin{lstlisting}
CREATE (:Empresa {eid: "emp_andina04", nombre: "Andina Logística",
                  industria: "Logística", tamano: "mediana", ...});

MATCH (a:Usuario {uid: "usr_valentina16"}), (b:Empresa {eid: "emp_andina04"})
CREATE (a)-[:TRABAJA_EN {cargo: "Directora de Tecnología",
                         fecha_inicio: "2021-03-01", actual: true}]->(b);
\end{lstlisting}

Guardamos las fechas como texto con el formato \cod{AAAA-MM-DD}, que es como
\emph{neomodel} las lee y escribe; con ese formato, el orden alfabético coincide con
el cronológico. Las contraseñas las guardamos cifradas, con el mismo método que usa Django.

Con los dos scripts de datos, la base queda con \textbf{175 nodos y 788 relaciones}.

\subsection{Cómo se ejecutan}

Incluimos en la aplicación un comando que ejecuta los tres scripts en orden:

\begin{lstlisting}[language={}]
python manage.py cargar_datos              # en una base vacía
python manage.py cargar_datos --limpiar    # borra todo y vuelve a cargar
\end{lstlisting}

También se pueden ejecutar a mano con \cod{cypher-shell}, en el mismo orden:
primero \cod{constraints}, luego \cod{seed_data} y al final \cod{datos_demo}.
Los pasos completos para instalar y correr el proyecto están en el archivo
\cod{README.md} del repositorio.

% =============================================================================
\capturafijatrue
\section{Consultas utilizadas}\label{sec:consultas}
% =============================================================================

Reunimos en \cod{queries.cypher} 17 consultas: las que responden las seis preguntas del
enunciado, las de caminos, las de recomendación y las que alimentan la analítica y el
grafo interactivo. Las consultas 1 a 14 tienen la misma lógica que las que ejecuta la
aplicación; la diferencia es que en el script devuelven columnas legibles, como nombres
y títulos, para poder probarlas en Neo4j Browser, mientras que la aplicación pide los
nodos completos. Las consultas 15 a 17 son de apoyo, para revisar la base. El archivo completo está en el Anexo~\ref{anx:consultas}.

Las palabras que empiezan con \cod{$}, como \cod{$uid}, son parámetros: la aplicación
pone ahí el usuario, la empresa o la oferta elegida en pantalla.

\begin{table}[htbp]
\centering
\caption{Preguntas del enunciado y la consulta que responde cada una.}
\rowcolors{2}{azulfila}{white}
\small
\begin{tabularx}{\linewidth}{@{}L c c@{}}
\cabecera \h{Pregunta} & \h{Consulta} & \h{Sección} \\
¿Qué usuarios tienen habilidades similares?                        & 3  & \ref{sec:similares} \\
¿Qué contactos en común existen entre dos profesionales?           & 2  & \ref{sec:comun} \\
¿Cuál es la ruta más corta entre un usuario y una empresa?         & 4  & \ref{sec:ruta} \\
¿Qué candidatos podrían recomendarse para una oferta?              & 6  & \ref{sec:candidatos} \\
¿Qué proyectos comparten tecnologías o miembros relacionados?      & 14 & \ref{sec:proyectos} \\
¿Qué habilidades son las más demandadas en las ofertas?            & 5  & \ref{sec:demandadas} \\
\bottomrule
\end{tabularx}
\end{table}

\subsection{Usuarios con habilidades similares}\label{sec:similares}

Partimos de una persona, pasamos a sus habilidades y de ahí a las demás personas que
tienen esas mismas habilidades. Cuantas más habilidades compartan, más arriba las ubicamos.

\begin{lstlisting}
MATCH (u1:Usuario {uid: $uid1})-[:TIENE_HABILIDAD]->(h:Habilidad)
      <-[:TIENE_HABILIDAD]-(u2:Usuario)
WHERE u1 <> u2
WITH u2, count(h) AS total_compartidas, collect(h.nombre) AS habilidades_comunes
ORDER BY total_compartidas DESC, u2.nombre ASC
RETURN u2.uid AS uid, u2.nombre + ' ' + u2.apellido AS nombre_completo,
       total_compartidas, habilidades_comunes
LIMIT $limit;
\end{lstlisting}

\captura[0.72\linewidth]{similares.png}{Personas con habilidades parecidas a las de Valentina Ríos.}{similares}

\subsection{Contactos en común}\label{sec:comun}

Un contacto en común es alguien conectado con las dos personas a la vez. El patrón
dibuja exactamente eso: una persona en el medio unida a ambas con \rel{CONECTA_CON}. Escribimos
las relaciones sin flecha porque el contacto es mutuo.

\begin{lstlisting}
MATCH (u1:Usuario {uid: $uid1})-[:CONECTA_CON]-(comun:Usuario)
      -[:CONECTA_CON]-(u2:Usuario {uid: $uid2})
WHERE u1 <> u2
RETURN DISTINCT comun.uid AS uid,
       comun.nombre + ' ' + comun.apellido AS nombre_completo,
       comun.titulo_profesional AS titulo
ORDER BY nombre_completo ASC;
\end{lstlisting}

\captura[0.6\linewidth]{comun.png}{Valentina Ríos y Laura Pineda tienen un contacto en común.}{comun}

\subsection{Ruta más corta entre un usuario y una empresa}\label{sec:ruta}

Usamos \cod{shortestPath} para buscar la cadena más corta, de hasta seis pasos, que une
a la persona con la empresa. Permitimos pasar por contactos (\rel{CONECTA_CON}) y por
empleos (\rel{TRABAJA_EN}), pero solo por empleos actuales: alguien que ya no trabaja en una
empresa no sirve como puente hacia ella.

\begin{lstlisting}
MATCH (u:Usuario {uid: $uid}), (e:Empresa {eid: $eid})
MATCH p = shortestPath((u)-[:CONECTA_CON|TRABAJA_EN*..6]-(e))
WHERE all(r IN relationships(p)
          WHERE type(r) <> 'TRABAJA_EN' OR coalesce(r.actual, true))
RETURN [n IN nodes(p) | {tipo: labels(n)[0], id: coalesce(n.uid, n.eid),
        nombre: CASE WHEN n:Usuario THEN n.nombre + ' ' + n.apellido
                     ELSE n.nombre END}] AS pasos_ruta,
       length(p) AS longitud_camino;
\end{lstlisting}

\captura[0.78\linewidth]{ruta_empresa.png}{Paula Gómez llega a Andina Logística en tres
pasos: conoce a Sofía Cárdenas, que conoce a Ricardo Salazar, que trabaja en la empresa.}{ruta-empresa}

Con la misma idea escribimos la consulta~1, entre dos personas y usando solo
\rel{CONECTA_CON}, y la consulta~13, que extiende el camino hasta una oferta con la
relación \rel{PUBLICA}
(Figura~\ref{fig:ruta-oferta}).

\captura[0.78\linewidth]{ruta_oferta.png}{Camino de Paula Gómez hasta una oferta de Andina
Logística: pasa por dos contactos y por la empresa, que publica la oferta.}{ruta-oferta}

\subsection{Candidatos para una oferta según conexiones y competencias}\label{sec:candidatos}

Para responder esta pregunta combinamos las dos cosas que pide. Primero buscamos a las
personas que cumplen al menos un requisito de la oferta (competencias). Después contamos
cuántos de sus contactos trabajan en la empresa que la publica (conexiones). Con eso
calculamos un puntaje:

\begin{center}
\small
puntaje $= 3 \times$ requisitos que cumple $+\; 2 \times$ requisitos obligatorios que cumple $+\; 2 \times$ contactos en la empresa
\end{center}

Dejamos fuera a quienes ya trabajan en esa empresa, porque no son candidatos.

\begin{lstlisting}
MATCH (o:Oferta {oid: $oid})-[r:REQUIERE]->(h:Habilidad)<-[:TIENE_HABILIDAD]-(u:Usuario)
WHERE NOT EXISTS {
  MATCH (u)-[t:TRABAJA_EN]->(:Empresa)-[:PUBLICA]->(o) WHERE coalesce(t.actual, true)
}
WITH o, u, count(DISTINCT h) AS total_coincidencias,
     collect(DISTINCT h.nombre) AS habilidades_match,
     count(DISTINCT CASE WHEN r.prioridad = 'obligatoria' THEN h END) AS obligatorias_cubiertas
OPTIONAL MATCH (u)-[:CONECTA_CON]-(c:Usuario)-[t:TRABAJA_EN]->(:Empresa)-[:PUBLICA]->(o)
WHERE coalesce(t.actual, true)
WITH o, u, total_coincidencias, habilidades_match, obligatorias_cubiertas,
     collect(DISTINCT c.nombre + ' ' + c.apellido) AS contactos
RETURN u.uid AS uid, u.nombre + ' ' + u.apellido AS nombre_completo,
       total_coincidencias, obligatorias_cubiertas, habilidades_match, contactos,
       total_coincidencias * 3 + obligatorias_cubiertas * 2 + size(contactos) * 2 AS score,
       EXISTS { (u)-[:POSTULA_A]->(o) } AS ya_postulado
ORDER BY score DESC, total_coincidencias DESC, nombre_completo ASC
LIMIT $limit;
\end{lstlisting}

\captura[0.72\linewidth]{candidatos.png}{Mejores candidatos para \emph{Ingeniero de Datos
Logísticos}. Debajo de cada nombre aparecen las habilidades que cumple y sus contactos
dentro de la empresa.}{candidatos}

\subsection{Proyectos que comparten tecnologías o miembros relacionados}\label{sec:proyectos}

Dado un proyecto, comparamos cada uno de los demás en tres aspectos y les damos un
puntaje:
\begin{itemize}
  \item \textbf{tecnologías en común} (2 puntos cada una), comparando las listas sin
        distinguir mayúsculas;
  \item \textbf{miembros en común}, personas que están en los dos equipos (3 puntos cada una);
  \item \textbf{miembros relacionados}, parejas de personas de cada equipo que están
        conectadas entre sí (1 punto cada pareja).
\end{itemize}

\begin{lstlisting}
MATCH (p:Proyecto {pid: $pid}), (otro:Proyecto)
WHERE otro <> p
WITH p, otro, [t IN coalesce(p.tecnologias, [])
               WHERE toLower(t) IN [x IN coalesce(otro.tecnologias, []) | toLower(x)]]
              AS tecnologias
OPTIONAL MATCH (p)<-[:PARTICIPA_EN]-(m:Usuario)-[:PARTICIPA_EN]->(otro)
WITH p, otro, tecnologias, collect(DISTINCT m) AS comunes
OPTIONAL MATCH (p)<-[:PARTICIPA_EN]-(a:Usuario)-[:CONECTA_CON]-(b:Usuario)
               -[:PARTICIPA_EN]->(otro)
WHERE NOT a IN comunes AND NOT b IN comunes
WITH otro, tecnologias, comunes,
     collect(DISTINCT a.nombre + ' ' + a.apellido + ' - ' + b.nombre + ' ' + b.apellido)
     AS conectados
WITH otro, tecnologias, conectados, [m IN comunes | m.nombre + ' ' + m.apellido]
     AS miembros_comunes,
     size(tecnologias) * 2 + size(comunes) * 3 + size(conectados) AS score
WHERE score > 0
RETURN otro.pid AS pid, otro.nombre AS nombre, tecnologias, miembros_comunes,
       conectados, score
ORDER BY score DESC, nombre ASC
LIMIT $limit;
\end{lstlisting}

\captura[0.8\linewidth]{proyectos.png}{Proyectos relacionados con \emph{Mapa de Rutas Andinas}.}{proyectos}

\subsection{Habilidades más demandadas}\label{sec:demandadas}

Contamos en cuántas ofertas abiertas aparece cada habilidad como requisito, separando
las obligatorias de las deseables. Las ofertas cerradas no las tenemos en cuenta.

\begin{lstlisting}
MATCH (h:Habilidad)<-[r:REQUIERE]-(o:Oferta {estado: 'abierta'})
RETURN h.hid AS hid, h.nombre AS habilidad, h.categoria AS categoria,
       count(o) AS total_ofertas,
       count(CASE WHEN r.prioridad = 'obligatoria' THEN 1 END) AS obligatorias,
       count(CASE WHEN r.prioridad = 'deseable' THEN 1 END) AS deseables
ORDER BY total_ofertas DESC, obligatorias DESC
LIMIT $limit;
\end{lstlisting}

\captura[0.72\linewidth]{demandadas.png}{Habilidades más pedidas en las ofertas abiertas.}{demandadas}

\subsection{Recomendaciones}

Con las consultas 7 a 10 generamos las sugerencias de la aplicación. Todas siguen la
misma idea: partimos de un nodo, exploramos lo que tiene cerca y sumamos puntos por cada
vínculo encontrado.

\begin{table}[htbp]
\centering
\caption{Consultas de recomendación y cómo puntúan.}
\rowcolors{2}{azulfila}{white}
\small
\begin{tabularx}{\linewidth}{@{}l c L@{}}
\cabecera \h{Recomienda} & \h{Consulta} & \h{Criterio} \\
Contactos                    & 7  & Personas sin conexión directa: 4 puntos por contacto en común, 3 por empresa en común y 2 por habilidad en común. \\
Empleos                      & 8  & Ofertas abiertas a las que la persona no se ha postulado, ordenadas por el porcentaje de requisitos que ya cumple. Muestra también los que le faltan. \\
Colaboradores para un proyecto & 9  & 5 puntos por cada habilidad que coincide con las tecnologías del proyecto y 3 por cada contacto en el equipo. \\
Talento para una empresa     & 10 & 4 puntos por cada habilidad pedida en sus ofertas abiertas y 3 por cada contacto que trabaja allí. \\
\bottomrule
\end{tabularx}
\end{table}

\subsection{Panel de analítica}\label{sec:analitica}

En el panel de analítica mostramos tres análisis de la red de contactos:
\begin{itemize}
  \item \textbf{Usuarios más influyentes}, con tres medidas: cuántos contactos tiene
        cada persona (centralidad de grado), su PageRank y su papel como puente entre
        grupos (intermediación). El PageRank considera influyente a quien está conectado
        con personas que, a su vez, son influyentes.
  \item \textbf{Comunidades}, detectadas con el algoritmo de Louvain, que agrupa a las
        personas que tienen muchas conexiones entre sí y pocas hacia afuera. A cada
        comunidad la nombramos con las habilidades más frecuentes de sus miembros.
  \item \textbf{Habilidades más demandadas}: en cuántas ofertas abiertas se pide cada
        habilidad, como en la consulta 5, comparado con cuántas personas de la red la
        tienen.
\end{itemize}

Neo4j entrega la red completa con la consulta 11, y sobre ella aplicamos los algoritmos
en Python:

\begin{lstlisting}
MATCH (a:Usuario)-[:CONECTA_CON]->(b:Usuario)
RETURN a.uid, b.uid;
\end{lstlisting}

\begin{nota}{¿Por qué los algoritmos corren en Python y no en Neo4j?}
Neo4j tiene una biblioteca para esto, \emph{Graph Data Science}, pero hay que instalarla
aparte. Preferimos que el proyecto funcione con una instalación normal de Neo4j en el
computador de cualquier integrante. Con 40 personas, el cálculo en Python tarda una
fracción de segundo. En una red grande convendría usar esa biblioteca, que hace el
cálculo dentro de la base sin sacar los datos.
\end{nota}

\captura[0.82\linewidth]{influencia.png}{Las siete personas más influyentes según su PageRank, comparado con su número de contactos.}{influencia}

\captura[0.8\linewidth]{comunidades.png}{Las cinco comunidades que encuentra el algoritmo de Louvain en la red de demostración.}{comunidades}

% =============================================================================
\capturafijafalse
\section{Investigación: ¿cuándo conviene un grafo y cuándo no?}\label{sec:investigacion}
% =============================================================================

\subsection{WheelsOn en un modelo relacional}

Para comparar, pensamos cómo se vería WheelsOn en tablas. Cada tipo de nodo sería una tabla
(\cod{usuario}, \cod{empresa}, \cod{oferta}…) y cada tipo de relación, una tabla
intermedia con las claves de sus dos extremos y sus propiedades:

\begin{lstlisting}[style=sql]
CREATE TABLE conexion (
  usuario_a TEXT REFERENCES usuario(uid),
  usuario_b TEXT REFERENCES usuario(uid),
  fecha     DATE,
  PRIMARY KEY (usuario_a, usuario_b)
);
\end{lstlisting}

Los dos modelos pueden guardar la misma información. Lo que cambia es el trabajo
que hace cada una para responder, y eso depende del tipo de pregunta.

\subsection{Dónde gana el grafo}

\subsubsection{Recorridos de profundidad variable}
La ruta más corta entre dos personas no tiene una cantidad fija de pasos: pueden ser
dos o seis. En Cypher son dos líneas (Sección~\ref{sec:ruta}). En SQL hace falta una
consulta recursiva que, nivel por nivel, extiende todos los caminos posibles y lleva la
cuenta de por dónde pasó cada uno para no dar vueltas:

\begin{lstlisting}[style=sql]
WITH RECURSIVE ruta(actual, camino, saltos) AS (
  SELECT 'usr_valentina16', ARRAY['usr_valentina16'], 0
  UNION ALL
  SELECT c.otro, r.camino || c.otro, r.saltos + 1
  FROM ruta r
  JOIN (SELECT usuario_a AS uno, usuario_b AS otro FROM conexion
        UNION ALL
        SELECT usuario_b, usuario_a FROM conexion) c ON c.uno = r.actual
  WHERE r.saltos < 6 AND NOT c.otro = ANY(r.camino)
)
SELECT camino, saltos FROM ruta
WHERE actual = 'usr_paula24'
ORDER BY saltos LIMIT 1;
\end{lstlisting}

Además de ser más larga, esta consulta hace más trabajo: en cada nivel vuelve a cruzar
la tabla de conexiones y explora todos los caminos de hasta seis pasos antes de elegir
el más corto. \cod{shortestPath} de Neo4j, en cambio, avanza desde los dos extremos a la
vez y se detiene en cuanto los caminos se encuentran.

La diferencia crece con la profundidad. En el libro \emph{Neo4j in Action} (2014) se comparan MySQL y Neo4j
buscando los amigos de los amigos en una red de un millón de personas con unos 50
contactos cada una:

\begin{table}[htbp]
\centering
\caption{Tiempo para encontrar contactos a distinta profundidad en una red de un millón de
personas, según el libro \emph{Neo4j in Action} (2014, cap.~1).}
\rowcolors{2}{azulfila}{white}
\begin{tabular}{@{}c r r r@{}}
\cabecera \h{Profundidad} & \h{MySQL (s)} & \h{Neo4j (s)} & \h{Personas encontradas} \\
2 & 0{,}016    & 0{,}010 & 2\,500 \\
3 & 30{,}267   & 0{,}168 & 110\,000 \\
4 & 1\,543{,}505 & 1{,}359 & 600\,000 \\
5 & no terminó & 2{,}132 & 800\,000 \\
\bottomrule
\end{tabular}
\end{table}

A dos pasos los dos responden al instante. Desde el tercero, cada \texttt{JOIN} más
multiplica el trabajo de la base relacional, mientras que Neo4j solo recorre los
contactos que encuentra. Las cifras exactas dependen del equipo y de cómo se
configure cada base, y quienes han intentado repetir el experimento no siempre han
obtenido una diferencia tan grande. Por eso lo tomamos como una ilustración de la
tendencia y no como una medida exacta.

\subsubsection{Preguntas que combinan varios tipos de relación}
{\sloppy La consulta de candidatos (Sección~\ref{sec:candidatos}) recorre cinco tipos de relación:
\rel{REQUIERE}, \rel{TIENE_HABILIDAD}, \rel{CONECTA_CON}, \rel{TRABAJA_EN} y
\rel{PUBLICA}. En Cypher el patrón se lee casi igual que la pregunta. En SQL serían al
menos seis tablas unidas con \texttt{JOIN}, y el sentido de la consulta se pierde entre
las condiciones de unión.\par}

\subsubsection{Un esquema que cambia}
Durante el desarrollo agregamos dieciséis tipos de relación, como las alianzas entre
empresas y todas las solicitudes pendientes. En Neo4j bastó con empezar a crearlas: las
consultas que ya existían no se vieron afectadas porque solo recorren los tipos que
nombran. En una base relacional, cada una habría sido una tabla nueva y una migración
de la base.

\subsubsection{Recomendaciones a partir de una persona}
Las recomendaciones parten de un nodo y exploran sus alrededores. Su costo depende del
tamaño de ese vecindario, no del tamaño de la red. Es el caso de uso más común de las
bases de grafos: sugerencias de contactos, detección de fraude o rutas logísticas.

\subsection{Dónde no gana el grafo}\label{sec:no-gana}

\subsubsection{Sumas y reportes sobre todos los datos}
``Las habilidades más demandadas'' no recorre la red: cuenta requisitos y los agrupa. En
SQL es un \texttt{GROUP BY} igual de simple que el Cypher de la Sección~\ref{sec:demandadas}:

\begin{lstlisting}[style=sql]
SELECT h.nombre, COUNT(*) AS ofertas
FROM requisito r
JOIN habilidad h ON h.hid = r.hid
JOIN oferta o    ON o.oid = r.oid
WHERE o.estado = 'abierta'
GROUP BY h.nombre
ORDER BY ofertas DESC;
\end{lstlisting}

Para este tipo de preguntas el grafo no aporta nada, y las bases relacionales tienen
décadas de optimización y herramientas de reportes alrededor.

\subsubsection{Reglas de integridad}
Una base relacional hace cumplir reglas por sí sola: claves foráneas, campos
obligatorios (\texttt{NOT NULL}) y condiciones (\texttt{CHECK}). La edición gratuita de
Neo4j solo garantiza que un valor no se repita. En WheelsOn, reglas como ``una oferta
pertenece a una sola empresa'', ``cada empresa tiene un solo administrador'' o ``no se
acepta a más personas que cupos'' tuvimos que programarlas nosotros y probarlas a
mano para asegurarnos de que los datos siguieran siendo consistentes.

\subsubsection{Datos con estructura fija y pocas relaciones}
Facturas, inventarios o registros contables tienen siempre los mismos campos y casi no
se recorren como red. En esos casos, las tablas son más simples, SQL es un estándar que
casi todo el mundo conoce y la base relacional es la mejor opción.

\subsubsection{Algoritmos sobre la red completa}
PageRank o la detección de comunidades no son recorridos cortos: procesan todos los
nodos varias veces. Ninguno de los dos modelos los resuelve con una consulta normal;
hace falta una biblioteca especializada o, como en WheelsOn, calcularlos fuera de la
base (Sección~\ref{sec:analitica}).

\subsubsection{Estandarización}
SQL es un estándar desde 1986 y casi todas las bases relacionales lo hablan. Los
lenguajes de grafos apenas se están unificando: el estándar ISO GQL, basado en buena
parte en Cypher, se publicó en 2024.

\subsection{Resumen}

\begin{table}[htbp]
\centering
\caption{Qué modelo conviene según el tipo de pregunta.}
\rowcolors{2}{azulfila}{white}
\small
\begin{tabularx}{\linewidth}{@{}L l L@{}}
\cabecera \h{Escenario} & \h{Conviene} & \h{Ejemplo en WheelsOn} \\
Caminos de longitud desconocida                     & Grafo      & Ruta más corta hacia una empresa. \\
Preguntas que cruzan muchos tipos de relación       & Grafo      & Candidatos por competencias y contactos. \\
Recomendaciones desde una persona o un proyecto     & Grafo      & Contactos, empleos y colaboradores sugeridos. \\
Esquema que cambia con frecuencia                   & Grafo      & Alianzas y solicitudes agregadas sobre la marcha. \\
Conteos y reportes sobre todo el conjunto           & Cualquiera & Habilidades más demandadas. \\
Reglas de integridad estrictas                      & Relacional & Cupos de una oferta y administrador único. \\
Datos tabulares con pocas relaciones                & Relacional & No hay en WheelsOn: sería el caso de pagos o facturación. \\
Algoritmos sobre la red completa                    & Biblioteca especializada & Influencia y comunidades. \\
\bottomrule
\end{tabularx}
\end{table}

Para elegir entre los dos hay que mirar qué tipo de preguntas se le van a hacer a los
datos. Si recorren relaciones a una profundidad que no se conoce de antemano, conviene
el grafo. Si suman, filtran o validan registros de
estructura fija, la base relacional es más simple y más segura. Muchos sistemas reales
usan las dos: una relacional para las transacciones y un grafo para las
recomendaciones.

% =============================================================================
\section{Conclusiones}
% =============================================================================

\begin{itemize}[label=\textcolor{azul}{\rule{0.5em}{0.5em}}, leftmargin=1.4em, itemsep=0.5em]
  \item Implementamos todas las funcionalidades del enunciado guardando la red completa
        en Neo4j, y respondimos sus seis preguntas con consultas Cypher que también usa
        la aplicación.
  \item Guardar en las relaciones los datos de cada vínculo, como el cargo, la etapa de
        una postulación o la prioridad de un requisito, fue la decisión de modelado que
        más simplificó las consultas.
  \item Separar lo pendiente de lo confirmado en tipos de relación distintos evitó que
        las rutas, las recomendaciones y la analítica tuvieran que filtrar estados en
        cada consulta.
  \item El grafo nos dio consultas cortas y legibles para caminos y recomendaciones,
        pero nos obligó a programar y probar a mano reglas de integridad que una base
        relacional habría hecho cumplir por sí sola.
  \item De la investigación concluimos que un grafo conviene cuando las
        preguntas recorren relaciones de profundidad variable, y una base relacional
        cuando se trata de sumar, filtrar y validar datos con estructura fija.
\end{itemize}

% =============================================================================
\section{Referencias}
% =============================================================================

\begin{itemize}[label=\textcolor{azul}{\rule{0.5em}{0.5em}}, leftmargin=1.4em, itemsep=0.4em]
  \item Neo4j, Inc. \emph{Cypher Manual}, versión 5. \url{https://neo4j.com/docs/cypher-manual/5/}
  \item Neo4j, Inc. \emph{Cypher Manual}: restricciones y su disponibilidad por edición. \url{https://neo4j.com/docs/cypher-manual/5/constraints/}
  \item ISO/IEC 39075:2024. \emph{Information technology --- Database languages --- GQL}.
  \item Page, L., Brin, S., Motwani, R. y Winograd, T. (1999). \emph{The PageRank Citation Ranking: Bringing Order to the Web}. Stanford InfoLab.
\end{itemize}

% =============================================================================
\appendix
% =============================================================================
\section{Script de inicio de la base de datos}\label{anx:inicio}
% =============================================================================

Este es el script completo con el que creamos las restricciones y los índices,
guardado en \cod{scripts/constraints.cypher}. Los dos scripts de datos suman cerca de
mil líneas de \cod{CREATE} con el mismo formato que mostramos en la
Sección~\ref{sec:scripts}, y están en la misma carpeta del repositorio.

\lstinputlisting[style=anexo]{../scripts/constraints.cypher}

% =============================================================================
\section{Script de consultas}\label{anx:consultas}
% =============================================================================

Este es el archivo completo con las consultas que usa la aplicación
(\cod{scripts/queries.cypher}).

\lstinputlisting[style=anexo]{../scripts/queries.cypher}

\end{document}
